% =====================================================================
%  Thème 4 — CORRIGÉ — Niveau DIFFICILE
% =====================================================================
\documentclass[11pt,a4paper]{article}
\usepackage[utf8]{inputenc}
\usepackage[T1]{fontenc}
\usepackage{lmodern}
\usepackage[french]{babel}
\usepackage{amsmath,amssymb,mathtools}
\usepackage{icomma}
\usepackage{xcolor}
\usepackage[most]{tcolorbox}
\usepackage{enumitem}
\usepackage{array}
\usepackage{booktabs}
\usepackage{fancyhdr}
\usepackage[hidelinks]{hyperref}
\usepackage[a4paper,margin=2.2cm,headheight=26pt,headsep=10pt]{geometry}

\definecolor{primary}{RGB}{223,87,99}
\definecolor{primarydark}{RGB}{150,42,55}
\definecolor{excol}{RGB}{0,135,90}

\tcbset{boxbase/.style={enhanced,breakable,boxrule=0.9pt,arc=2.5pt,
   left=3mm,right=3mm,top=2.3mm,bottom=2.3mm,
   fonttitle=\bfseries,coltitle=white,toptitle=0.6mm,bottomtitle=0.6mm}}
\newtcolorbox[auto counter]{correction}[1][]{boxbase,colback=excol!5,colframe=excol,
  title={\textsc{Corrigé}~\thetcbcounter\ifx\relax#1\relax\relax\else\textmd{ --- \itshape #1}\fi}}

\pagestyle{fancy}
\fancyhf{}
\fancyhead[L]{\small\textcolor{primarydark}{\textbf{Fiche 6} — Les limites}}
\fancyhead[R]{\small\textcolor{primarydark}{\textsc{Thème 4 — Corrigé Difficile}}}
\fancyfoot[C]{\small\thepage}
\renewcommand{\headrulewidth}{0.8pt}
\renewcommand{\headrule}{\hbox to\headwidth{\color{primary}\leaders\hrule height \headrulewidth\hfill}}

\newcommand{\R}{\mathbb{R}}

\begin{document}

\begin{center}
{\Large\bfseries\color{primarydark}Thème 4 — Radicaux : facteur dominant \& conjugué}\\[2pt]
{\large\color{excol}Corrigé — Niveau Difficile}
\end{center}
\vspace{6pt}

\begin{correction}[une étude aux deux infinis]
$f(x)=\sqrt{x^2-2x+5}-x$.
\begin{enumerate}[label=\alph*),leftmargin=2em,itemsep=3pt]
  \item $x^2-2x+5$ a pour discriminant $\Delta=4-20=-16<0$ : le trinôme est \emph{toujours} du signe
        de son coefficient dominant ($+1$), donc $>0$. La racine existe pour tout $x$ : $D_f=\R$.
  \item Conjugué $\sqrt{x^2-2x+5}+x$ :
  \[ f(x)=\frac{\bigl(\sqrt{x^2-2x+5}\bigr)^2-x^2}{\sqrt{x^2-2x+5}+x}=\frac{(x^2-2x+5)-x^2}{\sqrt{x^2-2x+5}+x}
     =\frac{-2x+5}{\sqrt{x^2-2x+5}+x}. \]
  \item En $+\infty$, $x>0$ donc $\sqrt{x^2-2x+5}=x\sqrt{1-\tfrac2x+\tfrac{5}{x^2}}$. On met $x$ en facteur :
  \[ f(x)=\frac{x\bigl(-2+\frac5x\bigr)}{x\bigl(\sqrt{1-\frac2x+\frac{5}{x^2}}+1\bigr)}
        =\frac{-2+\frac5x}{\sqrt{1-\frac2x+\frac{5}{x^2}}+1}\xrightarrow[x\to+\infty]{}\frac{-2}{1+1}=-1. \]
  Donc $\displaystyle\lim_{x\to+\infty} f(x)=-1$ (asymptote oblique $y=-x-1$ pour la courbe de $\sqrt{x^2-2x+5}$).
  \item En $-\infty$ : $\sqrt{x^2-2x+5}\to+\infty$ et $-x\to+\infty$ (car $x<0$). La somme
        $\sqrt{x^2-2x+5}-x=\sqrt{x^2-2x+5}+(-x)$ est donc $(+\infty)+(+\infty)$, \emph{sans indétermination} :
        \[ \lim_{x\to-\infty} f(x)=+\infty. \]
\end{enumerate}
\end{correction}

\begin{correction}[double conjugué]
$g(x)=\dfrac{\sqrt{x+1}-2}{\sqrt{x-2}-1}$.
\begin{enumerate}[label=\alph*),leftmargin=2em,itemsep=3pt]
  \item Conditions : $x+1\geq 0$, $x-2\geq 0$ et dénominateur non nul $\sqrt{x-2}\neq 1$ (soit $x\neq 3$).
        D'où $D_g=[2,3[\,\cup\,]3,+\infty[$.
  \item En $3$ : numérateur $\sqrt{4}-2=0$, dénominateur $\sqrt{1}-1=0$ : forme $\frac00$.
  \item On multiplie par les deux conjugués :
  \[ g(x)=\frac{(\sqrt{x+1}-2)(\sqrt{x+1}+2)\,(\sqrt{x-2}+1)}{(\sqrt{x-2}-1)(\sqrt{x-2}+1)\,(\sqrt{x+1}+2)}
        =\frac{(x+1-4)(\sqrt{x-2}+1)}{(x-2-1)(\sqrt{x+1}+2)}. \]
  Or $x+1-4=x-3$ et $x-2-1=x-3$, qui se simplifient :
  \[ g(x)=\frac{(x-3)(\sqrt{x-2}+1)}{(x-3)(\sqrt{x+1}+2)}=\frac{\sqrt{x-2}+1}{\sqrt{x+1}+2}. \]
  \item En $x=3$ : $\displaystyle\lim_{x\to 3} g(x)=\frac{\sqrt{1}+1}{\sqrt{4}+2}=\frac{2}{4}=\frac12$.
\end{enumerate}
\end{correction}

\begin{correction}[mise en situation : écart entre deux distances]
$e(x)=\sqrt{x^2+40x}-x$.
\begin{enumerate}[label=\alph*),leftmargin=2em,itemsep=3pt]
  \item $e(10)=\sqrt{100+400}-10=\sqrt{500}-10\approx 22{,}36-10=12{,}36$ ;\\
        $e(100)=\sqrt{10000+4000}-100=\sqrt{14000}-100\approx 118{,}32-100=18{,}32$.
        L'écart \emph{augmente} mais semble ralentir (il ne croît pas sans limite).
  \item Conjugué : $\displaystyle e(x)=\frac{(x^2+40x)-x^2}{\sqrt{x^2+40x}+x}=\frac{40x}{\sqrt{x^2+40x}+x}$.
        Pour $x>0$, $\sqrt{x^2+40x}=x\sqrt{1+\frac{40}{x}}$, donc
        \[ e(x)=\frac{40x}{x\bigl(\sqrt{1+\frac{40}{x}}+1\bigr)}=\frac{40}{\sqrt{1+\frac{40}{x}}+1}
        \xrightarrow[x\to+\infty]{}\frac{40}{1+1}=20. \]
  \item L'écart entre la distance par la route et la distance à vol d'oiseau se \textbf{stabilise vers
        $20$ km} pour les grands trajets : au-delà d'une certaine distance, allonger le trajet n'augmente
        pratiquement plus cet écart.
\end{enumerate}
\end{correction}

\end{document}
